\section{The three-point cosheaf construction}\label{app:fork-structure}

\subsection{Cosheaves and functor categories}

Write $O_0=\{c\}$, $O_1=\{a,c\}$, and $O_2=\{b,c\}$ for the nonempty
proper opens of $X$. Morphisms of the set-valued cosheaves defined in
Section~\ref{sec:cosheaves} are natural transformations.

\begin{proposition}\label{prop:fork-cosheaves}
Restriction to $O_0,O_1,O_2$ defines an equivalence
$\operatorname{CoSh}(X,\mathbf{Set})\simeq\mathbf{Set}^{J}$,
where $J=(1\leftarrow0\to2)$. A quasi-inverse sends a diagram
$Y_1\xleftarrow{r_1^Y}Y_0\xrightarrow{r_2^Y}Y_2$ to the cosheaf
with value $Y_j$ at $O_j$ and $Y_1\amalg_{Y_0}Y_2$ at $X$.
\end{proposition}
\begin{proof}
Let $F:\mathcal O(X)\to\mathbf{Set}$ satisfy $F(\varnothing)=\varnothing$.
Every cover of $O_j$ contains $O_j$. A cover of $X$ that omits $X$
contains both $O_1$ and $O_2$. Each of its members lies in $O_1$
or $O_2$, so a cocone on the cover is determined
by its maps on $O_1$ and $O_2$ agreeing on $O_0=O_1\cap O_2$.
Thus the remaining cosheaf condition is that the canonical map
\[
 F(O_1)\amalg_{F(O_0)}F(O_2)\longrightarrow F(X)
\]
be a bijection.

Extend $Y$ by sending $O_i\subseteq X$ to the coprojection
$\iota_i^Y:Y_i\to Y_1\amalg_{Y_0}Y_2$ for $i=1,2$. A morphism $f:Y\to Z$ of diagrams consists of functions
$f_j:Y_j\to Z_j$ satisfying $f_i r_i^Y=r_i^Z f_0$ for $i=1,2$.
It induces the unique map
\[
 f_X:Y_1\amalg_{Y_0}Y_2\longrightarrow Z_1\amalg_{Z_0}Z_2,
 \qquad f_X\iota_i^Y=\iota_i^Z f_i\quad(i=1,2).
\]
Uniqueness gives preservation of identities and composition.
Restriction of this extension recovers the original diagram and
its morphisms. For a cosheaf $F$, the comparison from the extension
of its restriction to $F$ is the identity on the proper opens and
the canonical bijection above at $X$. Naturality at $X$ follows
from the equations for $f_X$.
\end{proof}

For an object $Y$ of $\mathcal E$, the pushout has a
quotient presentation on the disjoint union $Y_1+Y_2$: identify
$(1,r_1^Yy)$ with $(2,r_2^Yy)$ for every $y\in Y_0$ and take the
equivalence relation $\sim_Y$ generated by these pairs. Enumerate the
generating pairs in both directions and the diagonal on $Y_1+Y_2$.
Finite paths in these edges enumerate $\sim_Y$.
For $f:Y\to Z$ in $\mathcal E$, the component at $X$ of the associated
cosheaf morphism is
\[
 (Y_1+Y_2)/{\sim_Y}\longrightarrow (Z_1+Z_2)/{\sim_Z},\qquad
 [(i,y)]\longmapsto [(i,f_i(y))]\quad(i=1,2).
\]
Indeed, $f_i r_i^Y=r_i^Z f_0$, so each generating pair maps to a
generating pair.

Products and subobjects are taken in $\mathcal E$. The cosheaf
associated with $Y\times Z$ has value
$(Y_1\times Z_1)\amalg_{Y_0\times Z_0}(Y_2\times Z_2)$ at $X$.
This need not be the product of the values at $X$ of
the cosheaves associated with $Y$ and $Z$. Write $0$, $1$, and $2$
for empty, singleton, and two-element numerical carriers, respectively.
For example,
the object $H=(1\leftarrow0\to1)$ of $\mathcal E$ satisfies
$H\times H\cong H$,
and the cosheaf associated with $H$ has two elements at $X$,
whereas the Cartesian square of that set has four.
Likewise $H\to(1\leftarrow1\to1)$ is a monomorphism in $\mathcal E$,
although the associated cosheaf morphism has component $2\to1$ at $X$.
We use regularity and regular images as in
Appendix~\ref{app:categorical-hypotheses}.
Finite coproducts are \emph{extensive} when, for every finite family
$(Y_k)_k$, copairing defines an equivalence
\[
 \prod_k(\mathcal E/Y_k)\longrightarrow
 \mathcal E/\Bigl(\coprod_kY_k\Bigr).
\]

\begin{proposition}\label{prop:fork-total}
The categories $\mathcal C_{\mathrm{rec}}$ and
$\mathcal E=\mathcal C_{\mathrm{rec}}^J$ have finite limits and
finite coproducts. The category $\mathcal E$ is regular and
extensive. Its monomorphisms are componentwise injections, and its
regular images are the componentwise c.e.\ images.
\end{proposition}
\begin{proof}
Proposition~\ref{prop:numerical-structure} gives finite limits,
disjoint pullback-stable finite coproducts, regular images, and
computable sections of surjections in $\mathcal C_{\mathrm{rec}}$.

Limits and finite coproducts in $\mathcal E$ are formed
componentwise. Their universal morphisms are natural because their
components are uniquely determined by the corresponding numerical
universal properties. For a natural transformation $f:Y\to Z$,
the equations $r_i^Zf_0=f_i r_i^Y$ show that its component images
form a subfunctor of $Z$.
For a componentwise surjection $q:Y\to Z$, a natural transformation
$u:Y\to W$ coequalizing its kernel pair descends componentwise.
If $v_jq_j=u_j$ are the descended morphisms, then
\[
 r_i^Wv_0q_0=r_i^Wu_0=u_i r_i^Y
 =v_iq_i r_i^Y=v_i r_i^Zq_0.
\]
Surjectivity of $q_0$ gives $r_i^Wv_0=v_i r_i^Z$, so the
descent is a natural transformation. These image factorizations are stable
under pullback because the corresponding assertion holds in
each component. They give regularity of $\mathcal E$.

The three objects of $\mathcal E$
\[
 H_0=(1\leftarrow1\to1),\qquad
 H_1=(1\leftarrow0\to0),\qquad
 H_2=(0\leftarrow0\to1)
\]
detect component elements: total morphisms $H_j\to Y$ are in bijection
with elements of $Y_j$. Each testing morphism $H_j\to Y$ has computable components because every
nonempty component of $H_j$ is a singleton.
If a morphism of $\mathcal E$ is monic,
testing it on these morphisms proves that each component is injective.
The converse follows by cancelling those component injections.
An injective morphism of $\mathcal C_{\mathrm{rec}}$ has a partial recursive inverse
on its c.e.\ image, obtained by preimage search.

Finally, for a morphism $Y\to A+B$, the inverse images of the two
coproduct summands form c.e.\ subfunctors $Y^A,Y^B$ and give an
isomorphism $Y^A+Y^B\cong Y$. Morphisms over $A+B$ restrict uniquely to morphisms
over $A$ and $B$, and two such morphisms combine uniquely by
copairing. This proves the slice equivalence expressing extensivity.
\end{proof}

\subsection{Partial maps, ranges, and connected restriction idempotents}

\begin{proposition}\label{prop:fork-structure}
The category $\mathcal A=\operatorname{Par}(\mathcal E,
\operatorname{Mono})$ is a split prodominical restriction category
with binary restriction products, restriction coproducts, weak
choice, and ranges satisfying the agreement law. The range of a morphism
represented by $Y\hookleftarrow D\xrightarrow{v}Z$ is the partial
identity on the regular image of $v$.
Joins of restriction idempotents are partial identities on finite
unions of their subfunctors, and inverse image along every partial
morphism preserves binary meets and finite joins. Ranges commute with products and
are stable under pullback along total morphisms.
\end{proposition}
\begin{proof}
Replace the monic leg of a representing span by its image
inclusion, using its computable inverse on that image. A partial
morphism $f:Y\rightharpoonup Z$ is then represented by a c.e.\ subfunctor
$D_f\subseteq Y$ and a natural transformation $D_f\to Z$.
For $j\in\{0,1,2\}$, write $f_j:Y_j\rightharpoonup Z_j$ for the
partial recursive function with domain $(D_f)_j$ whose restriction
to $(D_f)_j$ is the $j$-component of this natural transformation.
Use the same convention for other partial morphisms.
If $g:Z\rightharpoonup W$ has domain subfunctor $D_g$, the domain of $gf$
at component $j$ is
\[
 \{y\in(D_f)_j:f_j(y)\in(D_g)_j\},
\]
and its value there is $g_j(f_j(y))$. These carriers are c.e.\
and form a subfunctor by naturality. This is the
pullback composition of the representing spans. Associativity
and identities follow from these domain and value formulas.

The restriction $\bar f$ is the partial identity on $D_f$.
The restriction axioms reduce componentwise to the corresponding
identities of partial functions. Every restriction idempotent
splits through its domain subfunctor: the inclusion is total,
and its partial inverse is the identity with that domain.
Empty domains give restriction zeros. The restriction product
uses componentwise Cartesian products. For $f:W\rightharpoonup Y$
and $g:W\rightharpoonup Z$, the projection equations determine the pairing by
\[
 (D_{\langle f,g\rangle})_j=(D_f)_j\cap(D_g)_j,\qquad
 \langle f,g\rangle_j(w)=(f_j(w),g_j(w))
 \quad(w\in (D_f)_j\cap(D_g)_j).
\]
Coproducts use disjoint unions of carriers. The domain of a partial copairing is
the disjoint union of the two domains. A product with an empty domain
is empty, which proves the prodominical zero law.

The component images $f_j((D_f)_j)$ form a c.e.\ subfunctor
$I_f\subseteq Z$. Its partial identity satisfies $1_{I_f}f=f$.
If a restriction idempotent $e$ on $Z$ satisfies $ef=f$, its domain
contains $f_j((D_f)_j)$ at each $j$. Thus $1_{I_f}\leq e$, and
$\ran(f)=1_{I_f}$. For $g,h:Z\rightharpoonup W$ with
$gf=hf$, choose a preimage of each $z\in(I_f)_j$. Equality of the two partial composites shows
that $g_j$ is defined at $z$ exactly when $h_j$ is, and their
values there agree. Hence $g\ran(f)=h\ran(f)$, including
equality of domains.

For a total $u:W\to Z$, the projection
$D_f\times_Z W\to W$ has component image
$u_j^{-1}(f_j((D_f)_j))$. Thus its range is $u^*\ran(f)$,
which is the asserted pullback law. The image of a product
of two representing morphisms is the product of their images, giving
$\ran(f\times g)=\ran(f)\times\ran(g)$.
Intersections and finite unions of domain subfunctors are again
c.e.\ subfunctors; their partial identities give the meets and joins
of the corresponding restriction idempotents. For a restriction idempotent
$e$ on $Z$ with domain subfunctor $D_e$, the inverse image $f^*e$
has component carrier
\[
 \{y\in(D_f)_j:f_j(y)\in(D_e)_j\}.
\]
This formula proves preservation of binary meets and finite joins,
including the empty join.

Suppose that $p:Y\rightharpoonup Z$ is a partial monomorphism
and let $D_p$ be its domain subfunctor. If two elements of $(D_p)_j$
have the same image, their corresponding total morphisms
$H_j\to D_p\to Y$ have equal composites with $p$.
Partial monicity makes the original morphisms equal, since their
images lie in $D_p$. Thus each $p_j$ is injective on $(D_p)_j$.
Preimage search gives recursive inverse components on $I_p$;
they form a natural transformation by injectivity and naturality of
the representing morphism $D_p\to Z$.
The resulting partial morphism $p^\dagger:Z\rightharpoonup Y$ has
domain $I_p$ and satisfies $pp^\dagger=\overline{p^\dagger}$ and
$pp^\dagger p=p$.
This proves weak choice.
\end{proof}

\begin{proposition}\label{prop:fork-local}
The category $\mathcal A$ is locally connected.
\end{proposition}
\begin{proof}
Let $Y$ be an object of $\mathcal E$. For $y\in Y_0$, take the
finite subfunctor $D\subseteq Y$ with
\[
 D_0=\{y\},\qquad D_1=\{r_1^Yy\},\qquad D_2=\{r_2^Yy\}.
\]
The partial identity on $D$ is connected: in any partition of $D$
into two disjoint subfunctors, the part containing $y$ in component
$0$ must contain its images in components $1$ and $2$, and hence
must equal $D$. The partial identity on a singleton subfunctor
in component $i\in\{1,2\}$, with the other components empty,
is connected as well.

Suppose that partial morphisms $f,g:Y\rightharpoonup Z$ agree
after composition with every connected restriction idempotent.
Restriction to the
finite subfunctor generated by $y\in Y_0$ detects whether
$f_0$ and $g_0$ are defined at $y$, and compares their values
when defined. Partial identities on singleton subfunctors in
components $1$ and $2$ do the same for $f_1,g_1$ and $f_2,g_2$.
Thus all three component partial functions agree, proving $f=g$.
\end{proof}

\subsection{A universal object and a frame}

Let $V$ be the object of $\mathcal A$ defined in
\eqref{eq:universal-fork}.
\begin{lemma}\label{lem:fork-dominance}
Every object $Y$ of $\mathcal E$ admits a restriction monomorphism
$s_Y:Y\to V$ in $\mathcal A$.
\end{lemma}
\begin{proof}
For $j=0,1,2$, choose a recursive injection $c_j:Y_j\to\N$
with a partial recursive inverse on its image, using a tuple
coding of the ambient finite power. Define
\[
 s_{Y,i}(y_i)=(0,c_i(y_i))\quad(i=1,2),\qquad
 s_{Y,0}(y)=(0,c_0(y),c_1(r_1^Yy),c_2(r_2^Yy)).
\]
These maps are total and injective, and satisfy
$r_i^Vs_{Y,0}=s_{Y,i} r_i^Y$ for $i=1,2$.
Their images form a c.e.\ subfunctor $I\subseteq V$. The inverses
$s_{Y,j}^{-1}:I_j\to Y_j$ are computable and satisfy
$r_i^Y s_{Y,0}^{-1}=s_{Y,i}^{-1}r_i^I$. They define a partial morphism
$r:V\rightharpoonup Y$ with $rs_Y=1_Y$ and $s_Y r=\bar r$, the partial
identity on $I\subseteq V$.
\end{proof}

For $Y\in\mathcal E$, let $Y^\#$ be the object of nonempty words
defined in Section~\ref{sec:cosheaves}.
For a partial morphism $f:Y\rightharpoonup Z$ of $\mathcal A$
with domain subfunctor $D_f$, define
$f^\#:Y^\#\rightharpoonup Z^\#$ to have, at component $j$, domain
\[
 \{\operatorname{code}_{m(Y_j)}(y_1,\ldots,y_n):
          n\ge1,\ y_1,\ldots,y_n\in(D_f)_j\}\subseteq(Y^\#)_j.
\]
Its component at $j$ decodes a sequence, applies $f_j$ to every
entry, and encodes the resulting sequence in $(Z_j)^\#$.
These assignments preserve identities and composition and define
a functor $(-)^\#:\mathcal A\to\mathcal A$.
Concatenation and singleton insertion give morphisms
$\mu_Y:Y^\#\times Y^\#\to Y^\#$ and $\eta_Y:Y\to Y^\#$
in $\mathcal E$.
For an internal semigroup $(M,\mu)$ in $\mathcal E$, with
associative multiplication $\mu:M\times M\to M$, and a morphism
$f:Y\to M$, multiplication of the $f$-images of the entries gives
the unique semigroup morphism $\widehat f:Y^\#\to M$ satisfying
$\widehat f\eta_Y=f$.

Fix a recursive pairing bijection
$\pi:\N^2\to\N$. For $n\ge1$ choose tuple bijections
$\operatorname{encode}_n:\N^n\to\N$ and their inverses
$\operatorname{decode}_n:\N\to\N^n$, both computable uniformly in $n$.
Encode a nonempty word $(k_1,\ldots,k_n)$ by
$\pi(n-1,\operatorname{encode}_n(k_1,\ldots,k_n))$.

\begin{lemma}\label{lem:fork-frames}
There are isomorphisms in $\mathcal E$
$V\cong V\times V$, $V\cong V+V$, and $V\cong V^\#$.
\end{lemma}
\begin{proof}
For the product isomorphism, decode each of $k,z,x_1,x_2$ into a
pair using $\pi^{-1}$. Send $(k,z,x_1,x_2)$ to the resulting pair
of four-tuples. At component $i\in\{1,2\}$, decode $(k,x)$ into
the resulting pair of two-tuples. The naturality equations hold
coordinatewise, and coordinatewise pairing gives the inverse.

For the coproduct isomorphism, write $k=2q+t$ with
$t\in\{0,1\}$. Send a tuple with first coordinate $k$ to
summand $t$, replace $k$ by $q$, and retain its remaining
coordinates. This rule is the same at all three components.
The inverse replaces $q$ in summand $t$ by $2q+t$.

For the word isomorphism, decode $k$ as the nonempty word
$(k_1,\ldots,k_n)$. Apply $\operatorname{decode}_n$ to $z,x_1,x_2$
to obtain $(z_j)_j,(x_{1,j})_j,(x_{2,j})_j$, respectively. Send $(k,z,x_1,x_2)$ to the numerical code
\[
 \operatorname{code}_4
 ((k_1,z_1,x_{1,1},x_{2,1}),\ldots,(k_n,z_n,x_{1,n},x_{2,n}))
 \in(V^\#)_0.
\]
At endpoint $i$, the input is $(k,x_i)$. The same decodings give
\[
 \operatorname{code}_2((k_1,x_{i,1}),\ldots,(k_n,x_{i,n}))\in(V^\#)_i.
\]
These maps form a natural transformation. To invert them, decode
the sequence, recover the word of first coordinates, and apply
the tuple encodings to the remaining coordinate sequences.
These inverses are total recursive.
\end{proof}

\subsection{A Turing morphism}\label{app:fork-turing}

We construct a Turing morphism $\tau:V\times V\rightharpoonup V$
in $\mathcal A$ by giving a monomorphism $D_\tau\hookrightarrow V\times V$
and a natural transformation $D_\tau\to V$.

We use an acceptable numbering
$\varphi:\N\times\N\rightharpoonup\N$ of the partial recursive
functions and write $\varphi_p(t)=\varphi(p,t)$.
The specialization property of this numbering is the following:
if $m\geq0$, $n\geq1$, and
$G:\N^m\times\N^n\rightharpoonup\N$ is partial recursive,
there is a total recursive map $s:\N^m\to\N$ such that
\[
 \varphi_{s(y)}(\operatorname{encode}_n x)\simeq G(y,x).
\]
Here $\simeq$ means equality of partial values,
including simultaneous undefinedness. The case $m=0$ is a
constant program index. This is the $s$-$m$-$n$ theorem.

\begin{theorem}\label{thm:fork-turing}
The object $V$ of
$\mathcal A=\operatorname{Par}(\mathcal C_{\mathrm{rec}}^J,
\operatorname{Mono})$ admits a Turing morphism
$\tau:V\times V\rightharpoonup V$: for every object $Y$ of
$\mathcal A$ and every morphism $F:Y\times V\rightharpoonup V$
in $\mathcal A$, there is a total morphism
$p_F:Y\to V$ such that $\tau(p_F\times1_V)=F$.
At component $i\in\{1,2\}$, this Turing morphism is given by
\begin{equation}\label{eq:fork-turing-endpoint}
 \tau_i((k,p),(\ell,x))
 =\operatorname{decode}_2
     \varphi_p(\operatorname{encode}_2(\ell,x)).
\end{equation}
\end{theorem}

\begin{proof}
For $i=1,2$, let $D_{\tau,i}\subseteq V_i\times V_i$ be the c.e.\ set
on which the numerical computation in
\eqref{eq:fork-turing-endpoint} halts, and define
$\tau_i:V_i\times V_i\rightharpoonup V_i$ with domain $D_{\tau,i}$
by that equation. The first coordinate $k$
of the parameter has no effect on this computation.

At component $0$, write the parameter as
$P=(k,p_0,p_1,p_2)\in V_0$ and the input as
$x=(\ell,x_0,x_1,x_2)\in V_0$.
Define partial recursive maps $\psi_j:V_0\times V_0\rightharpoonup V_j$
for $j=0,1,2$ by the following parallel computations:
\begin{align*}
 \psi_0(P,x)&\simeq
   \operatorname{decode}_4
      \varphi_{p_0}(\operatorname{encode}_4 x),\\
 \psi_i(P,x)&\simeq
   \operatorname{decode}_2
      \varphi_{p_i}(\operatorname{encode}_2(\ell,x_i))
      \qquad(i=1,2).
\end{align*}
Let $D_{\tau,0}\subseteq V_0\times V_0$ consist of those pairs $(P,x)$
for which all three computations halt and
\begin{equation}\label{eq:fork-turing-check}
 \psi_i(P,x)=r_i^V \psi_0(P,x)\qquad(i=1,2).
\end{equation}
Define $\tau_0:V_0\times V_0\rightharpoonup V_0$ with domain
$D_{\tau,0}$ by $\tau_0(P,x)=\psi_0(P,x)$.
Dovetailing the three computations and checking equality of their
numerical outputs enumerates $D_{\tau,0}$ and computes $\tau_0$ on it.

If $(P,x)\in D_{\tau,0}$, then
$(r_i^V P,r_i^V x)=((k,p_i),(\ell,x_i))\in D_{\tau,i}$.
Thus $D_\tau=(D_{\tau,0},D_{\tau,1},D_{\tau,2})$ is a c.e.\ subfunctor of $V\times V$.
Equation~\eqref{eq:fork-turing-check} gives
\[
 r_i^V\tau_0(P,x)
   =\tau_i(r_i^V P,r_i^V x)
 \qquad((P,x)\in D_{\tau,0}).
\]
The restrictions $\tau_j|_{D_{\tau,j}}:D_{\tau,j}\to V_j$, for $j=0,1,2$,
therefore form a morphism $D_\tau\to V$ in $\mathcal E$.
Together with the inclusion $D_\tau\hookrightarrow V\times V$,
they define the required partial morphism $\tau$.

Now let $F:Y\times V\rightharpoonup V$ be any partial morphism.
Choose its representation by a c.e.\ subfunctor
$D\subseteq Y\times V$ and a morphism $D\to V$ in
$\mathcal E$. Its components are the restrictions
$F_j|_{D_j}:D_j\to V_j$, where $j=0,1,2$.
Let $d_j=m(V_j)$, so $d_0=4$ and $d_1=d_2=2$.
For each component define a partial recursive numerical function
\[
 G_j:\N^{m(Y_j)}\times\N^{d_j}\rightharpoonup\N
\]
as follows. On input $(y,x)$, wait until $(y,x)$ appears in an
enumeration of $D_j$; then compute $F_j(y,x)$ and return its
$\operatorname{encode}_{d_j}$-code. This function has domain
exactly $D_j$.
Specialization gives total recursive maps
$s_j:\N^{m(Y_j)}\to\N$ satisfying
\begin{equation}\label{eq:fork-specialization}
 \operatorname{decode}_{d_j}
   \varphi_{s_j(y)}(\operatorname{encode}_{d_j}x)
 \simeq F_j(y,x),
 \qquad\text{with domain }D_j.
\end{equation}

Define the component maps of $p_F:Y\to V$ by
\begin{align*}
 p_{F,i}(y_i)&=(0,s_i(y_i)) &&(i=1,2),\\
 p_{F,0}(y)&=(0,s_0(y),s_1(r_1^Yy),s_2(r_2^Yy)).
\end{align*}
These maps are recursive and total on their respective carriers:
the maps $s_j$ are total on the ambient numerical powers and the
maps $r_i^Y$ are total on $Y_0$. Moreover,
$r_i^Vp_{F,0}(y)=(0,s_i(r_i^Yy))=p_{F,i}(r_i^Yy)$,
so $p_F$ is a natural transformation in $\mathcal E$, hence total in
$\mathcal A$.

At component $i\in\{1,2\}$, equations
\eqref{eq:fork-turing-endpoint} and
\eqref{eq:fork-specialization} show that the partial function
\[
 (y_i,v_i)\longmapsto\tau_i(p_{F,i}(y_i),v_i)
 :Y_i\times V_i\rightharpoonup V_i
\]
has domain $D_i$ and agrees there with $F_i$.
At component $0$, fix $(y,x)\in Y_0\times V_0$.
If $(y,x)\in D_0$, the computation $\psi_0(p_{F,0}(y),x)$ returns $F_0(y,x)$.
Since $D$ is a subfunctor,
$(r_i^Yy,r_i^Vx)\in D_i$ for $i=1,2$.
The computations $\psi_i(p_{F,0}(y),x)$, for $i=1,2$, consequently
halt, and naturality of the morphism $D\to V$ gives
\[
 F_i(r_i^Yy,r_i^Vx)=r_i^VF_0(y,x).
\]
Both tests in \eqref{eq:fork-turing-check} succeed, so
$\tau_0(p_{F,0}(y),x)=F_0(y,x)$.
Conversely, if $\tau_0(p_{F,0}(y),x)$ is defined, the computation
$\psi_0(p_{F,0}(y),x)$ has halted. Equation~\eqref{eq:fork-specialization}
then implies $(y,x)\in D_0$.
Thus $\tau(p_F\times1_V)=F$, with equality of domains and values
at each component.
\end{proof}

\subsection{The recursion category and its splitting}

\begin{proposition}\label{prop:fork-recursion}
Let $\mathscr R_X$ be the full subcategory of $\mathcal A$
whose objects are obtained from $V$ by finitely many applications
of the chosen binary products, binary coproducts, and nonempty-word formation.
Then $\mathscr R_X$ is a recursion category. Its inclusion
into $\mathcal A$ is a range functor, and its restriction
splitting is equivalent to $\mathcal A$.
\end{proposition}
\begin{proof}
The isomorphisms of Lemma~\ref{lem:fork-frames}, applied inductively,
show that every object of $\mathscr R_X$ is isomorphic to $V$.
The binary restriction products, coproducts, and prodominical
zero law restrict from $\mathcal A$ to this full subcategory.
The Turing morphism of Theorem~\ref{thm:fork-turing} has
source $V\times V$ and target $V$ in the subcategory.
Its total parameter morphisms for objects of $\mathscr R_X$ belong
to $\mathscr R_X$ by fullness. Hence it retains its Turing
universal property there, proving that $\mathscr R_X$ is a
recursion category.

By fullness, the restriction idempotents on each object of $\mathscr R_X$
are the same as in $\mathcal A$. Hence the least $e$ satisfying $ef=f$
is the same in both categories. The inclusion preserves restrictions,
ranges and the chosen coproducts, so it is a range functor.

For a restriction idempotent $e$ on an object $Y$ of $\mathscr R_X$,
write $\iota_e:D_e\to Y$ and $r_e:Y\rightharpoonup D_e$ for its
domain inclusion and partial inverse. Send $(Y,e)$ to $D_e$
and a morphism $f:(Y,e)\rightharpoonup(Z,d)$, represented by
$f:Y\rightharpoonup Z$ satisfying $f=dfe$, to
$r_d f \iota_e:D_e\rightharpoonup D_d$.
The equations $\iota_er_e=e$ and $r_e\iota_e=1_{D_e}$ show that
identities and composition are preserved. Conversely, send a
partial map $g:D_e\rightharpoonup D_d$ to $\iota_d g r_e$.
This morphism belongs to $\mathscr R_X$ by fullness and satisfies
$d(\iota_d g r_e)e=\iota_d g r_e$. The two assignments are inverse on
each hom-set, so the functor is fully faithful.
It preserves restrictions, since
$\overline{r_d f \iota_e}=r_e\bar f \iota_e$.
For every $Z\in\mathcal A$, Lemma~\ref{lem:fork-dominance}
provides a restriction monomorphism $s_Z:Z\to V$ with partial inverse
$s_Z^\dagger:V\rightharpoonup Z$. The domain subfunctor of
$s_Zs_Z^\dagger$ is isomorphic to $Z$,
so the functor is essentially surjective.
\end{proof}

\begin{corollary}\label{cor:fork-karoubi}
Every idempotent in $\mathcal A$ splits, and the inclusion
$\mathscr R_X\hookrightarrow\mathcal A$ extends, up to natural
isomorphism, to a restriction equivalence
$Q:\operatorname{Kar}(\mathscr R_X)\to\mathcal A$.
\end{corollary}
\begin{proof}
Let $p:Y\rightharpoonup Y$ satisfy $p^2=p$, and let
$D\subseteq Y$ be its domain subfunctor. For $j=0,1,2$, put
$I_j=p_j(D_j)\subseteq Y_j$. These sets are c.e., and naturality
of $p:D\to Y$ gives $r_i^Y(I_0)\subseteq I_i$ for $i=1,2$.
Thus $I\subseteq Y$ is an object of $\mathcal E$.
Let $\iota_I:I\to Y$ be the inclusion, and let
$r:Y\rightharpoonup I$ have domain $D$ and component functions $p_j$.
The equation $p^2=p$ implies that $I_j\subseteq D_j$ and that
$p_j$ fixes every element of $I_j$. Hence $\iota_I r=p$ and $r\iota_I=1_I$.

For each idempotent $e:Y\rightharpoonup Y$ in $\mathscr R_X$, write
$I_e$ for its image subfunctor and
$\iota_e:I_e\to Y$, $r_e:Y\rightharpoonup I_e$ for the splitting just
constructed. Define
$Q:\operatorname{Kar}(\mathscr R_X)\to\mathcal A$ by
$Q(Y,e)=I_e$ and
$Q(f)=r_d f \iota_e:I_e\rightharpoonup I_d$ for
$f:(Y,e)\rightharpoonup(Z,d)$.
The equations $\iota_er_e=e$ and $r_e\iota_e=1_{I_e}$ prove the functor
laws. As in Proposition~\ref{prop:fork-recursion}, the inverse
assignment on hom-sets is $g\mapsto \iota_d g r_e$, so $Q$ is fully faithful.
For $f:(Y,e)\rightharpoonup(Z,d)$, restriction preservation follows from
$\overline{r_d f \iota_e}=r_e\bar f e \iota_e$, since $e$ fixes its image
and $f=dfe$ takes values in the image of $d$.
Every object of $\mathcal A$ is isomorphic to the image of a restriction
idempotent on $V$, so this functor is essentially surjective.
\end{proof}

Use the restriction equivalence
\[
 \mathsf Q_{\mathrm{num}}:\Split(\mathrm{Rec})
       \longrightarrow\mathcal K_{\mathrm{rec}}
       =\operatorname{Par}(\mathcal C_{\mathrm{rec}},\operatorname{Mono})
\]
constructed in Proposition~\ref{prop:numerical-structure}.
Sending a total morphism of $\mathcal K_{\mathrm{rec}}$ to its underlying
function gives the isomorphism
$\Tot(\mathcal K_{\mathrm{rec}})\cong\mathcal C_{\mathrm{rec}}$.
Under this identification, restriction to total categories gives the
equivalence $F:\mathcal B_{\mathrm{rec}}\to\mathcal C_{\mathrm{rec}}$
of that proposition.

Decidable equality on $V$ has the meaning fixed in
Appendix~\ref{app:categorical-hypotheses}: it is a morphism
$\epsilon:V\times V\to\Two$ in $\mathcal E$ whose inverse image
of $\mathbf t:1\to\Two$ is the diagonal subfunctor.

\begin{proposition}\label{prop:fork-distinction}
The recursion category $\mathscr R_X$ is not equivalent, as a
restriction category, to the recursion category of partial recursive functions.
The object $V$ has no decidable equality in $\mathcal E$.
\end{proposition}
\begin{proof}
The terminal object of $\mathcal E$ is the constant functor at a singleton.
A subobject with empty component $0$ has independently empty or
singleton components $1$ and $2$, giving four subobjects. A
subobject with singleton component $0$ must also have singleton
components $1$ and $2$, giving one more. Thus its subobject order has five elements, whereas
the terminal singleton of $\mathcal C_{\mathrm{rec}}$ has two.
An equivalence of the asserted restriction categories would induce
an equivalence of their restriction splittings, hence of their
total categories $\mathcal E$ and $\mathcal C_{\mathrm{rec}}$.
Equivalences preserve terminal objects and their subobject orders,
so this is impossible.

Suppose that $\epsilon$ decides equality on $V$. The functor $1+1$
is constant and sends both nonidentity morphisms of $J$ to
identity functions. Choose distinct $u,v\in V_0$ differing only
in their $z$ coordinate. Then $r_1^Vu=r_1^Vv$, so
$\epsilon_1(r_1^Vu,r_1^Vv)=\mathbf t_1(*)$.
Naturality gives $\epsilon_0(u,v)=\mathbf t_0(*)$, contrary to $u\ne v$.
\end{proof}


